\section{Robustness under spatially inhomogeneous couplings}
\label{sec:robustness}

We now extend the single homogeneous certificate to a class of perturbed
energies. The extension preserves a positive \emph{global} annihilation
barrier and gives a uniform dynamical stability statement. It does not
identify the perturbed minimax with the height of a particular continued
saddle.

Write $m_0,s_+,s_-$ for the certified unperturbed minimum and saddles,
and put
\[
 h=\E(s_+)=\E(s_-),\qquad B_0=h-\E(m_0).
\]
We use the certified domain $\mathcal D$, whose closure misses $\K$, and the
certified upper path $\Gamma_0$ from $m_0$ to $\K$, so that
\begin{equation}
 \label{eq:robust-seed-global}
 m_0\in\mathcal D,\qquad \E|_{\partial\mathcal D}\ge h,
 \qquad \max_t\E(\Gamma_0(t))=h,
 \qquad .00567867<B_0<.00570116.
\end{equation}
All operator norms and gradients in this section refer to the Riesz
identification associated with the $\X$ inner product.

For $b\in\ell^1(\mathcal E;\R)$ define
\begin{equation}
 \label{eq:coupling-energy}
 E_b(u)=\sum_{e\in\mathcal E}(1+b_e)(1-\cos D_eu)
       =\E(u)+W_b(u).
\end{equation}
The endpoint set remains $\K$ when $1+b_e>0$. Denote by $B_b(a)$ the
annihilation barrier for $E_b$ starting from $a$: its paths are all strong
$\X$-continuous paths from $a$ to $\K$, with no prescribed charge history,
symmetry, or exterior configuration.

\begin{lemma}[Uniform differential and oscillation bounds]
\label{lem:coupling-bounds}
Let $p=\nrm{b}_{\ell^1}$. Then $W_b$ is smooth on $\X$, and
\begin{align}
 \nrm{\nabla_{\X}W_b(u)}&\le\frac{p}{\sqrt2},&
 \nrm{H_{W_b}(u)}&\le\frac p2,\label{eq:coupling-derivatives}\\
 \nrm{H_b(u)}&\le1+\frac p2,&
 \nrm{H_b(u)-H_b(v)}&\le
       \frac{1+p/2}{\sqrt2}\nrm{u-v}_{\X}.
 \label{eq:coupling-hessian}
\end{align}
Moreover, the absolutely convergent sums
\begin{equation}
 \label{eq:coupling-oscillation}
 W_-=2\sum_{b_e<0}b_e,\qquad W_+=2\sum_{b_e>0}b_e
\end{equation}
satisfy $W_-\le W_b(u)\le W_+$ for every $u\in\X$ and
$W_+-W_-=2p$.
\end{lemma}

\begin{proof}
Let $g_e$ be the Riesz representative of $D_e$. The nearest-neighbour
resistance identity gives $\nrm{g_e}_{\X}=1/\sqrt2$. Hence
\[
 \nabla_{\X}W_b(u)=\sum_e b_e\sin(D_eu)g_e,\qquad
 H_{W_b}(u)=\sum_e b_e\cos(D_eu)g_e\otimes g_e.
\]
The series converge absolutely in $\X$ and in operator norm. Their
derivatives converge locally uniformly; for derivatives of positive order
this also follows directly by summing $|b_e|\nrm{g_e}_{\X}^{j}$.
This proves smoothness and \eqref{eq:coupling-derivatives}. The unperturbed
bounds are $\nrm{H_0(u)}\le1$ and
$\nrm{H_0(u)-H_0(v)}\le\nrm{u-v}_{\X}/\sqrt2$. For the perturbation,
the Lipschitz bound on the cosine gives
\[
 \nrm{H_{W_b}(u)-H_{W_b}(v)}
 \le\sum_e|b_e|\nrm{g_e}_{\X}^3\nrm{u-v}_{\X}
 =\frac{p}{2\sqrt2}\nrm{u-v}_{\X}.
\]
This yields \eqref{eq:coupling-hessian}. Finally
for each bond the sign of $b_e$ gives
\[
 2\min(b_e,0)\le b_e(1-\cos D_eu)\le2\max(b_e,0).
\]
Absolute summability permits adding these inequalities over all bonds.
They hold for every compatible phase field, regardless of dependencies
between its increments. Thus $W_-$ and $W_+$ are lower and upper
enclosures; they need not be attainable extrema, and no independent
variation of the bond angles is required.
\end{proof}

The stationary-state certificate provides the following safe rational
information:
\begin{equation}
 \label{eq:robust-seed-gap}
 H_0(m_0)\succeq\sigma I,\qquad
 \nrm{H_0(s_\pm)^{-1}}\le\sigma^{-1},\qquad
 \operatorname{ind}H_0(s_\pm)=1,
 \qquad \sigma=\frac{57}{1000}.
\end{equation}
Indeed, the minimum's candidate coercivity exceeds $.07567$ and its
correction radius is less than $.017422$, while the saddle's candidate
inverse gap exceeds $.069925$ and its correction radius is less than
$.012825$. Subtracting the correction radii using the conservative
Hessian Lipschitz constant $1$ gives $.058248$ and $.057100$, respectively,
both greater than $\sigma$. The minimum's global principal-increment
branch margin is greater than $11/10$.

\begin{theorem}[A quantitative class of robust global barriers]
\label{thm:robustness}
Let
\[
 \delta=\frac1{500},\qquad r=\frac1{32},\qquad
 \tau=\frac{71}{100},\qquad g=\frac{541}{16000}.
\]
For every $b\in\ell^1(\mathcal E;\R)$ with $\nrm{b}_{\ell^1}\le\delta$,
the following assertions hold.
\begin{enumerate}
\item The energy $E_b$ has unique critical points $m_b,s_b^+,s_b^-$ in
the closed radius-$r$ balls about $m_0,s_+,s_-$, respectively.
These branches depend smoothly on $b$. The minimum satisfies
$H_b(m_b)\succeq gI$; each saddle has index one and inverse norm at most
$g^{-1}$. The two saddle branches remain distinct.
\item The minimum $m_b$ has precisely the original charges $+1$ and $-1$
at $(0,0)$ and $(12,0)$ and belongs to $\mathcal D$.
\item Its global annihilation barrier satisfies
\begin{equation}
 \label{eq:robust-barrier}
 \big|B_b(m_b)-B_0\big|\le2\nrm{b}_{\ell^1}.
\end{equation}
In particular, throughout this coupling class,
\begin{equation}
 \label{eq:robust-barrier-numeric}
 .00167867<B_b(m_b)<.00970116.
\end{equation}
\end{enumerate}
\end{theorem}

\begin{proof}
Since $\tau^2>1/2$, it is a rational upper bound for $1/\sqrt2$.
Fix an unperturbed critical point $z$ among the three certified states.
Taylor expansion gives
\[
 \nabla_{\X}\E(z+v)=H_0(z)v+N_z(v),\qquad
 \nrm{N_z(v)}\le\frac\tau2\nrm{v}_{\X}^2.
\]
On the radius-$r$ ball, $N_z$ has Lipschitz constant at most $\tau r$.
The equation $\nabla_{\X}E_b(z+v)=0$ is equivalent to the fixed point
equation
\[
 v=T_{b,z}(v):=-H_0(z)^{-1}
             \big[N_z(v)+\nabla_{\X}W_b(z+v)\big].
\]
By \eqref{eq:robust-seed-gap} and
Lemma~\ref{lem:coupling-bounds}, the exact self-mapping and contraction
tests are
\begin{align}
 \tau\delta+\frac\tau2r^2
 &=\frac{45227}{25600000}
 <\frac{45600}{25600000}=\sigma r,
 \label{eq:robust-newton-map}\\
 \operatorname{Lip}(T_{b,z})
 &\le\frac{\tau r+\delta/2}{\sigma}
 =\frac{371}{912}<1.
 \label{eq:robust-newton-contraction}
\end{align}
Banach's theorem produces the unique critical point $z_b$ in this ball.
At that point,
\begin{equation}
 \label{eq:robust-operator-perturbation}
 \nrm{H_b(z_b)-H_0(z)}\le\tau r+\frac\delta2
 =\frac{371}{16000},\qquad
 \sigma-\frac{371}{16000}=g.
\end{equation}
For the minimum this gives the asserted coercivity. For a saddle,
\eqref{eq:robust-seed-gap} implies that the spectrum of its self-adjoint
Hessian misses $(-\sigma,\sigma)$. Its negative spectral space is
one-dimensional and its complementary spectral space is positive.
The perturbation in \eqref{eq:robust-operator-perturbation} leaves the
quadratic form negative on the former space and positive on the latter.
Thus its index is at least one and at most one. The inverse exists by
the Neumann-series test, with norm at most
$1/(\sigma-371/16000)=g^{-1}$. This argument does not identify a
numerical saddle direction with an exact eigenvector. Smooth dependence
on the $\ell^1$ parameter follows from the smooth implicit-function
theorem and uniqueness; the strict self-mapping test keeps each root
in the interior of its ball.

For completeness, distinctness uses the recorded reaction-coordinate
enclosures. If $Q:\R^2\to\X$ is the reaction embedding and $G=Q^*Q$,
the certificate gives $\lambda_{\min}(G)>1/4$. Therefore the linear
functional consisting of the second coordinate of $G^{-1}Q^*u$ has
norm less than $2$. The candidate saddle coordinates differ by more
than $.89$ in this coordinate, and each original saddle correction has
norm less than $.013$. The exact coordinate difference is consequently
greater than $.89-4(.013)=.838$, and the original saddle distance is
greater than $.419>.4$. The continued saddle distance is then greater
than $.4-2r=27/80$.

Every bond increment changes by at most $\tau r$ between $m_0$ and
$m_b$. Its distance from the principal-branch cut remains greater than
\[
 \frac{11}{10}-\tau r=\frac{3449}{3200}>0.
\]
The straight segment crosses no branch cut, so all plaquette charges
are preserved. To establish membership in $\mathcal D$, use stationarity
of $m_0$ and the global bound $\nrm{H_0}\le1$:
\begin{equation}
 \label{eq:robust-low-bridge}
 \E\big(m_0+t(m_b-m_0)\big)
 \le\E(m_0)+\frac{r^2}{2}
 =\E(m_0)+\frac1{2048}<h,
 \qquad 0\le t\le1.
\end{equation}
Since the segment starts in $\mathcal D$, it cannot meet its boundary,
where \eqref{eq:robust-seed-global} gives $\E\ge h$. In particular
$m_b\in\mathcal D$. No inclusion of a full norm ball in $\mathcal D$
has been assumed.

It remains to prove the global estimate with the shifted origin. On
the radius-$r$ ball about $m_0$, both energies are strictly convex:
\[
 H_0\succeq(\sigma-\tau r)I>0,\qquad H_b\succeq gI>0.
\]
Stationarity at the respective minima and convexity on this ball imply
\begin{equation}
 \label{eq:robust-endpoint-energies}
 \E(m_b)\ge\E(m_0),\qquad E_b(m_b)\le E_b(m_0),\qquad
 \max_{[m_b,m_0]}E_b\le E_b(m_0).
\end{equation}
Every admissible path from $m_b$ to $\K$ must cross
$\partial\mathcal D$. Its maximal perturbed energy is at least
$h+W_-$, whereas
$E_b(m_b)\le E_b(m_0)\le\E(m_0)+W_+$. This proves
\[
 B_b(m_b)\ge B_0-(W_+-W_-).
\]
For the opposite inequality, prepend the bridge from $m_b$ to $m_0$
to $\Gamma_0$. Its energy is at most $h+W_+$: this follows on
$\Gamma_0$ from \eqref{eq:robust-seed-global}, and on the bridge
from \eqref{eq:robust-endpoint-energies}. Finally
$E_b(m_b)\ge\E(m_0)+W_-$, again by
\eqref{eq:robust-endpoint-energies}. Hence
$B_b(m_b)\le B_0+(W_+-W_-)$, which proves
\eqref{eq:robust-barrier} and \eqref{eq:robust-barrier-numeric}.
\end{proof}

\begin{corollary}[Uniform local stability for the energy-space flow]
\label{cor:flow}
For every coupling field in Theorem~\ref{thm:robustness}, the flow
\begin{equation}
 \label{eq:robust-gradient-flow}
 \dot u=-\nabla_{\X}E_b(u)
\end{equation}
is globally well posed in $\X$. If $\nrm{u(0)-m_b}_{\X}\le1/100$, then
the solution stays in that ball and
\begin{equation}
 \label{eq:robust-exponential-stability}
 \nrm{u(t)-m_b}_{\X}
 \le e^{-13t/500}\nrm{u(0)-m_b}_{\X},\qquad t\ge0.
\end{equation}
Along every solution,
$\frac{d}{dt}E_b(u)=-\nrm{\nabla_{\X}E_b(u)}_{\X}^2$.
\end{corollary}

\begin{proof}
The gradient is globally Lipschitz, with constant at most
$1+\delta/2=1001/1000$, by \eqref{eq:coupling-hessian}; the Hilbert-space
ODE therefore has a unique global solution. Moreover
\[
 \frac{1+\delta/2}{\sqrt2}<\frac{711}{1000}.
\]
Within distance $1/100$ of $m_b$, its Hessian is bounded below by
\[
 \left(\frac{541}{16000}-\frac{711}{100000}\right)I
 =\frac{10681}{400000}I\succ\frac{13}{500}I.
\]
For $w=u-m_b$, stationarity and integration along the segment from
$m_b$ to $u$ give
\[
 \frac{d}{dt}\frac{\nrm{w}_{\X}^2}{2}
 =-\inner{w}{\nabla_{\X}E_b(m_b+w)-\nabla_{\X}E_b(m_b)}
 \le-\frac{13}{500}\nrm{w}_{\X}^2.
\]
A first-exit argument makes the closed ball invariant; Gronwall's
inequality gives \eqref{eq:robust-exponential-stability}. The energy
identity follows by the chain rule.
\end{proof}

\begin{remark}[Compatibility with the effective exterior estimate]
The class above also fits the Green estimate of
Section~\ref{sec:green}, Theorem~\ref{thm:green}. For a continued critical
state $z_b$, stationarity gives
\[
 Dz_b=P_{\mathrm{grad}}q_b,\qquad
 q_b=Dz_b-(1+b)\sin(Dz_b),\qquad
 \nrm{q_b}_{\ell^1}\le\frac{\nrm{z_b}_{\X}^3}{6}
                         +\nrm{b}_{\ell^1}.
\]
Beyond the certified unperturbed exterior radius, the original increments
are bounded by $.018$. Continuation therefore gives
$|D_ez_b|\le .018+\tau/32=643/16000<.041$. On those exterior bonds,
\[
 |q_{b,e}|\le\left(\frac{.041^2}{6}+\nrm{b}_{\ell^\infty}\right)
                  |D_ez_b|
 \le\frac{13681}{6000000}|D_ez_b|,
 \qquad 2\frac{13681}{6000000}<1.
\]
Thus the required exterior source term remains absorbable uniformly
throughout the $\ell^1$ ball. In particular, a translated coupling
perturbation is not an independent additive far-field force: it retains
the factor $\sin(D_ez_b)$. No compact-support condition on $b$ is needed
for this implication.
\end{remark}

\begin{remark}[Scope of the robustness statement]
The budget $\nrm{b}_{\ell^1}\le1/500$ is a \emph{total} coupling budget.
For one modified bond it allows a relative change of $0.2\%$; it does
not allow that change independently on every bond of a growing box.
Nonzero constant rescalings of the whole lattice are outside this
$\ell^1$ class. The theorem proves neither that one of $s_b^{\pm}$
attains the perturbed minimax nor that reflection breaking selects a
different escape mechanism. Those conclusions require additional
separator and connection checks. Corollary~\ref{cor:flow} concerns the
specified $\X$ mobility, not the local $\ell^2$ gradient flow or a thermal
escape rate. The continuation and oscillation arguments are standard;
their role here is to turn the reproduced global certificate into an
explicit class of robust infinite-lattice conclusions.
\end{remark}
